\documentclass[12pt, titlepage]{article}
%\usepackage[nynorsk]{babel}
\usepackage[utf8]{inputenc}
\usepackage[T1]{fontenc}
\usepackage{minted, enumerate}
\usepackage[colorlinks = true,
linkcolor = blue,
urlcolor = blue]{hyperref}


\begin{document}

\title{Innlevering: Mappe}
\author{IIRA2001 - Haust 2023}
\date{}
\maketitle

Mappen består av dei to mappeoppgåvene vi har jobba med i haust og eit refleksjonsnotat

\section*{Del I}

I første del av mappen skulle de skrive 3 pythonprogram:
\begin{itemize}
   \item {Eitt som brukar den logistiske likninga, til dømes ved populasjonsvekst}
   \item {Ein sparekalkulator}
   \item {Og ein lånekalkulator}
\end{itemize}
Sjå eigen oppgåveomtale på blackboard for meir detaljar om programma

~\\
\textit{\textbf{MERK:} De blei òg bedt om å skrive eit refleksjonsnotat i del I. Dette skal ikkje leverast.
De leverer i staden \textbf{eitt} refleksjonsnotat som gjeld heile semesteret}

\section*{Del II}

I andre del av mappen skulle de lage 2 program som analyserte datasett, brukte eit web-api eller begge deler.
Hovudsakleg ville vi ha 2 program som gjorde noko interessant med 1 eller fleire datasett, der det eine programmet brukte eit web-api.
Sjå oppgåveomtale på blackboard for detaljar om programma og unntak til omtalen over

\section*{Skildring av programa}

\begin{itemize}
   \item{Til kvart pythonprogram skulle de òg skrive litt tekst. }
   \item {Teksten må innehalde ei skildring av kva målet med programmet er, til dømes kva problemstilling du undersøkjer.}
   \item {I tillegg treng du ein lita diskusjon og/eller skildring av resultatet av programmet. }
\end{itemize}

Kort sagt treng vi ein liten rapport som skildrar kva målet med koden din var, og diskuterer kva resultatet av koden din blei, eller kva du fann ut.\\
Kor lang teksten er, avhenger av programmet ditt. 
Dersom du har laga ein sparekalkulator og bruker han til å finne ut kor mykje du må spare per månad for å kjøpe deg ein lamborghini i 30-årsgåve til deg sjølv, er det kanskje færre vurderingar og diskusjonar å skrive enn om du undersøkjer effekten straumprisen i Europa har på aksjekursen til norsk hydro.

~\\
De kan levere kode og «rapport» på 2 måtar:
\begin{itemize}
   \item {Som 1 *.ipynb fil (Jupyter-notebook fil) per program med kodeceller for python og tekstceller for skildring/resultat/diskusjon}
   \item {Som 1 *.py fil (pythonscript) med koden din og 1 PDF-fil med skildring/resultat/diskusjon, per program}
\end{itemize}

\color{blue}
\subsubsection*{ENDRING}
Etter å ha sett på løysingane dykkar og gitt tilbakemeldingar, har eg merkt meg at mange har andre fornuftige måtar å organisere mappen sin på.
Dersom de vil lage eitt dokument med refleksjonsnotat,  skildring og diskusjon av programma går det òg fint.
Nokre har presentert mappen sin med streamlit og andre meir avanserte måtar - dette er òg greitt, men ver for all del forsiktig med å sørge for at sensoren klarer å køyre programmet/applikasjonen dykkar.

De har fridom til å presentere mappen dykkar på litt andre måtar enn slik det er forskrive over, men på eigen risiko:
\begin{itemize}
   \item {Vi må kunne køyre koden dykkar utan for mykje styr}
   \item {Koden bør vere skilt ut i eigne filar}
   \item {Programma treng tekst som skildrar dei, og diskuterer eventuelle problemstillingar og resultat}
\end{itemize}

\color{black}


\section*{Refleksjonsnotat}

Det siste som skal leverast er eit refleksjonsnotat i ein PDF-fil.

~\\
I løpet av semesteret har vi sett på grunnleggjande programmering i python, og anvendelse av dette i dataanalyse og økonomiske fag.
Siste del av mappen består av eit \textit{refleksjonsnotat} der de reflekterer rundt læringsprosessen de har vore gjennom, og korleis programmering kan vere fagleg relevant i vidare studie og verksemd.

Eit refleksjonsnotat er ein personleg tekst der du reflekterer over eigne tankar, erfaringar eller innsikt knytt til ein bestemt hendelse eller prosess (slik som læringsprosessen de no har vore gjennom).
Ved å bearbeide og reflektere rundt desse tankane og erfaringane er målet å fremje læring og personleg og fagleg vekst.

~\\
I hovudtrekk skal de reflektere rundt korleis programmering kan vere fagleg relevant, kva ferdigheiter de har tileigna dykk og korleis de ser føre dykk framtidig utvikling og bruk av det de har lært.\\
%\textbf{Tips, relevante punktar å ha med, og relevant spørsmål å reflektere rundt:}
\subsubsection*{Tips, relevante punktar og spørsmål}
\begin{itemize}
   \item{Legg vekt på refleksjon, analytisk og kritisk tenking}
   \item{Pek ut dei viktigaste programmeringsferdighetene de har utvikla}
   \item{Korleis har programmering styrka forståinga av økonomiske eller andre faglege aspekt?}
   \item{Har nokon av oppgåvene eller casene bidratt til forbetra forståing av økonomiske problemstillingar?}
   \item{Kva utfordringar møtte du, og korleis handterte du desse?}
   \item{Har du lært noko som er relevant for ein karriere du ser føre deg?}
   \item{De kan diskutere potensielle anvendelsar}
   \item{Identifiser nokre ting som var vellykka, og kvifor}
   \item{Korleis kan du byggje på det du har lært i framtida?}
   \item{Er det kanskje noko i oppgåvene de har levert de kunne ønskje å forbetre?}
   \item{Diskuter og grei ut om kva de har lært}
   \item{De kan reflektere over eiga tilnærming til å lære programmering}
   \item{Det vil vere lurt å knyte refleksjonar til konkrete døme i mappen}
   \item{Refleksjonsnotat er ein personleg tekst, så ver ærleg og open}
   \item{Det kan være lurt med ein oppsummering og ein slags konklusjon på slutten}
\end{itemize}

De treng naturlegvis ikkje å oppfylle alle punkta over, men dei kan vere til hjelp for å kome i gong.

~\\
\textit{Refleksjonsnotatet bør vere på rundt 750-1500 ord}

\section*{Inspera}

Det opnar for innlevering på Inspera mellom 18.12 og 22.12.
Heile mappen med kode og tekst pakkast saman til ei zip-fil og lastast opp i inspera.

Vi legg ut litt info og lenker på blackboard om korleis de pakkar saman filer til eit zip-arkiv, og kor de finn Jupyter-notebook filene dykkar.


\end{document}
